\chapter*{Introduction to This Volume: How Life Preserves and Rewrites Information}
\addcontentsline{toc}{chapter}{Introduction to This Volume}
\markboth{Introduction to This Volume}{}

The first two volumes answer how matter makes up life: elements form molecules, molecules form cells, and a metabolic network operates inside each cell. But life does something else even more remarkable: it \textbf{remembers}. A fertilized egg ``remembers'' how to develop into a complete body; you ``remember'' your parents' appearance; a species ``remembers'' the marks left by hundreds of millions of years of environmental change.

These memories are written in the same kind of molecule: DNA. Volume III tells how life preserves, copies, regulates, and rewrites information. Part VIII begins with pea flowers, follows Mendel as he turns inheritance into calculable rules, and traces how scientists established DNA as the hereditary material. Part IX enters the molecular machinery: how billions of letters are copied, transcribed, and translated into proteins; how genes are switched on, turned down, or turned up; and why ``one gene determines one trait'' is almost always wrong. Part X stretches time across hundreds of millions of years, asking how mutation, selection, and drift write evolutionary history. Part XI returns to the present: how sequencing, gene editing, genetic testing, and genetic modification are changing medicine, agriculture, and each of us---and which questions still have no answers.

Bring this volume's most important attitude with you from the start: genes matter, but genes are not destiny. Traits emerge from genes, environment, developmental history, and random processes together. Understanding that protects you more than memorizing terminology. It protects you from genetic fortune-telling and from being frightened by extreme stories of genetic omnipotence or genetic horror.
